\documentclass[12pt]{article}
% Préambule commun à tous les documents extraits de « La Revue CAPES »
\usepackage[T1]{fontenc}
\usepackage[utf8]{inputenc}
\usepackage{lmodern}
\usepackage{amsmath,amssymb,amsthm,amsfonts,mathrsfs}
\usepackage{setspace}
\usepackage{listings}
\usepackage{pgfplots}
\pgfplotsset{compat=1.18}
\usepackage{longtable}
\usepackage{adjustbox}
\usepackage{lettrine}
\usepackage{tkz-tab}
\usepackage{changepage}
\usepackage{pdflscape}
\usepackage{graphicx}
\usepackage{tikz}
\usepackage{xcolor}
\usepackage[a4paper,top=2cm,bottom=2.5cm,left=2.5cm,right=2cm]{geometry}
\usepackage{fancyhdr}
\usepackage{draftwatermark}
\SetWatermarkText{La RevueCapes}
\SetWatermarkLightness{0.9}
\SetWatermarkScale{2}
\usepackage{hyperref}
\hypersetup{colorlinks=true,linkcolor=blue!50!black,urlcolor=blue!50!black}

\definecolor{revue}{RGB}{20,60,120}

\pagestyle{fancy}
\fancyhf{}
\renewcommand{\headrulewidth}{0.4pt}
\setlength{\headheight}{15pt}

% \entete{Matière}{Titre}{Type de document}
\newcommand{\entete}[3]{%
  \fancyhead[L]{\small\textsc{La Revue CAPES} -- #1}%
  \fancyhead[R]{\small #2}%
  \fancyfoot[C]{\thepage}%
  \thispagestyle{plain}%
  \begin{center}
    {\color{revue}\rule{\textwidth}{1.2pt}}\\[4pt]
    {\small\textsc{Concours directs CAP-CEG / CAPES -- Burkina Faso}}\\[6pt]
    {\LARGE\bfseries\color{revue} #2}\\[6pt]
    {\large\itshape #1 -- #3}\\[2pt]
    {\color{revue}\rule{\textwidth}{1.2pt}}
  \end{center}
  \vspace{0.5cm}%
}

% Encadré pour les remarques ajoutées dans les corrigés
\newcommand{\remarque}[1]{\par\medskip\noindent\fcolorbox{revue}{revue!6}{\parbox{\dimexpr\linewidth-2\fboxsep-2\fboxrule}{\textbf{\color{revue}Remarque.} #1}}\par\medskip}


\begin{document}
\entete{Mathématiques}{Sujet type n°11}{Énoncé}
\begin{spacing}{1.5}

\textbf{Exercice 1}

1. Calculer les intégrales généralisées suivantes :

a . $\int_{0}^{+\infty}\frac{1}{(1+e^t)(1+e^{-t})}dt$ \hspace{1cm} b. $\int_{1}^{+\infty}\frac{e^{-\sqrt{t}}}{\sqrt{t}}dt$.

\medskip c. $\int_{0}^{+\infty}\frac{\arctan t}{1+t^2}dt$ \hspace{1cm}d. $\int_{1}^{+\infty}\frac{lnt}{t^2}dt$ \hspace{1cm}e. $\int_{a}^{+\infty}\frac{dt}{t(t+b)},a>0,b>0$



2. Étudier la nature des intégrales suivantes :

a) $\int_0^{+\infty}\frac{lnt}{t^2+1}dt$ \hspace{1cm}b) $\int_0^{+\infty}\frac{1-\cos t}{t^2}dt$ \hspace{1cm}c) $\int_0^{+\infty}\frac{\sin t+\cos t}{\sqrt{t^3+2}}dt$ \hspace{1cm}d) $\int_{-\infty}^{+\infty}\sin t dt$.

3. Étudier, suivant les valeurs de $\alpha\in\mathbb{R}$, la convergence des intégrales généralisées suivantes :

a. \[\int_0^{+\infty}\frac{dt}{t^{\alpha}}dt\]


b. \[\int_0^{+\infty}\frac{t-\sin t}{t^{\alpha}}dt\]


\textbf{Exercice 2}

 La ville de Ouagadougou a acquis un parc de photocopieurs auprès du ministre de l'éducation. On estime que la probabilité qu'un photocopieur tombe en panne au cours d'une journée est $p=0,002$. Ce parc de photocopieurs fait l'objet d'un contrat de maintenance prévoyant le remplacement du matériel si la durée de la panne excède la journée. On admettra, pour simplifier, que la machine en panne est toujours remise en service au début du jour ouvrable qui suit celui de la panne mais n'est jamais remise en service dans la journée de la panne.

1) Calculer, en justifiant, la probabilité, pour une machine donnée, de tomber en panne au moins une fois au cours d'une période de 40 jours.

2) Soit $X$, la variable aléatoire "nombre de pannes survenant au cours d'une période de 40 jours sur un photocopieur donné". Quelle loi de probabilité suit
 la variable aléatoire $X$ ? Justifier votre réponse.
 
 3) Calculer la probabilité, pour une machine donnée, de tomber en panne au moins une fois au cours d'une période de 40 jours, en utilisant le résultat de la question précédente.
 
 4) Le parc acheté est composé de 52 machines qui sont utilisées dans des conditions identiques et ont une probabilité p'=0,077 de tomber au moins une
 fois en panne au cours d'une période de 40 jours. On désignera $Y$, la variable aléatoire "nombre de photocopieurs du parc tombant en panne au moins une
 fois au cours d'une période de 40 jours". Calculer l'espérance mathématique de $Y$, sa variance et son écart-type.

\medskip \textbf{Exercice 3 }

 Soit $f$ la fonction  définie sur $\mathbb{R}^2$ par :
 
 $
 f(x,y) = \frac{2x^3+3xy^2}{x^2+y^2}$ si $(x,y)\neq (0,0)$ et $f(0,0) = 0$
 
 1. Montrer que $f$ est continue sur $\mathbb{R}^2$.

 2. Calculer les dérivées directionnelles en $(0,0)$ dans toutes les directions, puis calculer $\overrightarrow{grad}f(0,0)$.
 
 3. $f$ est-elle différentiable en $(0,0)$?
 
 \medskip\textbf{Exercice 4}

On pose $f(x)=\int_{0}^{1}\frac{t^{x-1}}{1+t}dt$

1. Déterminer le domaine de définition de $f$.

2. Montrer que $f$ est continue sur son domaine de définition.

3. Calculer $f(x)+f(x+1)$ pour tout $x>0$

4. En déduire un équivalent de $f$ en 0.

5. Déterminer la limite de $f$ 

\quad

\quad

\end{spacing}
\end{document}
